\documentclass[conference]{IEEEtran}

% ---- Packages ----
\usepackage[utf8]{inputenc}       % UTF-8 encoding
\usepackage{cite}                 % Sorted & compressed numeric citations [1]-[3]
\usepackage{amsmath}              % Math environments
\usepackage{graphicx}             % Include images: \includegraphics{}
\usepackage{url}                  % Line-breakable URLs
\usepackage[colorlinks=true,
            allcolors=black]{hyperref}  % Clickable links, but black text
\usepackage{caption}
\usepackage{subcaption}
% ---- Title & Authors ----
\title{Efficient Federated Split Learning: Synergizing Activation Quantization and Adaptive Synchronization}

\author{\IEEEauthorblockN{Baoyi Liu, Guanghua Liu, Jiale Liu, Mingyuan Shao, Wenjie Ding, Yi Sin Lin, Zhuolun Li}
\IEEEauthorblockA{School of Computing, Newcastle University, Newcastle upon Tyne, United Kingdom \\
Email: \{c5016854, c5003419, c5053016, c5011444, c5053577, c4032640, c5036432\}@newcastle.ac.uk}}

% ==================================================================
\begin{document}
\maketitle

% ---- Abstract ----
\begin{abstract}
  
\end{abstract}

% ---- Keywords (optional) ----
\begin{IEEEkeywords}
  1
\end{IEEEkeywords}

% ==================================================================
% SECTIONS — each teammate can edit their own section
% ==================================================================

\section{Introduction}
\label{sec:intro}

xxxxxxxxx


\section{Related Work}
\label{sec:system}

2. Related Work
This section reviews existing studies related to communication-efficient distributed learning systems for IoT environments. Existing research can be broadly categorised into three directions:
(1) communication-efficient Federated / Split Learning,
(2) adaptive scheduling and system-level optimisation, and
(3) machine learning methods for rainfall prediction in resource-constrained environments.
 
2.1 Communication-efficient Federated / Split Learning
Communication-efficient learning has become an important research topic in distributed machine learning, particularly in IoT and edge environments where bandwidth is limited. In traditional Federated Learning (FL), communication cost is mainly caused by frequent gradient or parameter exchange between clients and the server. To reduce transmission overhead, many studies focus on compression-based methods. For example, compressive sensing–based approaches dynamically adjust compression ratios to reduce bidirectional communication cost [5], while gradient quantisation and sparsification techniques are designed to address bandwidth constraints in wireless edge learning scenarios [6]. These methods show clear improvements in communication efficiency for conventional FL systems.
However, the communication pattern in Federated Split Learning (FSL) differs significantly from standard FL. Instead of exchanging full model gradients, clients send intermediate activations (also called smashed activations) to the server during the forward pass, and gradients are returned during the backward pass. As neural networks become deeper, the size of these activations can increase substantially, leading to high communication overhead. Therefore, compression techniques designed for traditional FL do not always directly address the communication characteristics of FSL.
Recent research attempts to improve efficiency in Split Learning and Federated Split Learning through architectural and synchronisation optimisation. Dynamic split point selection allows the training workload to be adjusted between client and server to reduce communication cost under different conditions [4]. Asynchronous training mechanisms relax strict synchronisation requirements to reduce communication bottlenecks [8]. Other works redesign the learning architecture to reduce server-side state dependency and improve scalability [9;10], while adaptive aggregation strategies aim to balance communication cost and model convergence in distributed environments [11]. Although these methods improve efficiency, they often assume relatively stable network conditions or introduce additional system complexity.
 
2.2 Adaptive Scheduling and System-Level Optimisation
In addition to compression and architectural changes, another research direction focuses on adaptive scheduling and system-level optimisation for distributed learning. These approaches aim to improve performance by considering device capability, network latency, and resource availability during training.
Computation-oriented optimisation methods attempt to reduce the workload on resource-limited client devices. Techniques such as model pruning, adaptive layer reduction, and lightweight network design are widely used in edge learning scenarios. Adaptive pruning strategies dynamically adjust model complexity according to device capability to reduce computation cost [7]. Lightweight prediction models are also designed to balance accuracy and efficiency in resource-constrained environments, particularly for time-series prediction tasks such as rainfall forecasting [2].
However, these approaches mainly optimise computational cost rather than communication overhead. Even with smaller models, clients still participate in frequent synchronisation rounds, which can lead to high cumulative communication cost in distributed systems. In real IoT environments, network latency, bandwidth variation, and device heterogeneity may further increase training time and reduce system stability. Therefore, latency-aware training, resource-aware scheduling, and adaptive communication control are increasingly considered necessary for practical distributed learning systems.

2.3 Machine Learning for Rainfall Prediction in IoT Environments 

Machine learning has been widely applied to rainfall prediction and environmental sensing, especially with the increasing availability of IoT sensor data. Lightweight prediction models are often used to improve efficiency while maintaining acceptable accuracy in resource-constrained environments [2]. These methods demonstrate that machine learning can provide useful prediction capability for real-time monitoring systems. 

However, rainfall prediction remains a challenging task due to data imbalance, temporal variability, and the difficulty of predicting extreme events. Many studies focus primarily on improving modelling accuracy, without fully considering deployment constraints such as communication cost, device capability, and distributed training requirements. In real IoT sensing networks, prediction systems must operate under limited bandwidth, unstable connections, and heterogeneous devices, which makes communication efficiency as important as model accuracy. 

Across the three research directions above, a common limitation can be observed. Most existing studies optimise only one aspect of the system, such as compression, model structure, or distributed architecture, without jointly considering data representation precision, activation size, and synchronisation frequency. As a result, the interaction between communication cost, model convergence, and system scalability remains insufficiently explored, particularly for rainfall prediction systems deployed in real IoT environments. 

To address these limitations, this work proposes a communication-efficient Federated Split Learning framework that jointly optimises activation compression and synchronisation frequency. By integrating spatial compression, temporal communication control, and adaptive scheduling mechanisms, the proposed approach aims to improve the practicality of rainfall prediction systems in real-world IoT sensing networks. 
\subsection{System Design}

xxxxxxxx
\subsubsection{Activation Compression (Spatial Optimization)}


\section{High-level design}
\label{sec:design}
xxxxxxxxxxx
\subsection{System Overview}

xxxxxxxxxxx


\begin{enumerate}
xxxxxxxx
\end{enumerate}




\subsection{Design Rationale}




% ---- References ----
% Compile order: pdflatex -> bibtex -> pdflatex -> pdflatex

\bibliographystyle{IEEEtran}
\bibliography{IEEEabrv,refs}


\appendix
\section{Team Contribution Table}

\begin{center}
\begin{tabular}{|p{3.5cm}|p{3cm}|p{9cm}|}
\hline
\textbf{Team Member} & \textbf{Role} & \textbf{Contribution} \\
\hline

Guanghua Liu & Introduction Writing &
Reviewed paper \cite{10522472}. Wrote the introduction section and explained the motivation of using Federated Split Learning for rainfall prediction. \\
\hline

Baoyi Liu & Literature Review &
Reviewed paper \cite{10478872}. Read related papers and summarized the main ideas for the related work section. \\
\hline

Yi Sin Lin & System Design &
Reviewed paper \cite{10078005}. Helped design the overall system architecture and described the client--server structure. \\
\hline

Mingyuan Shao & Activation Compression &
Reviewed paper \cite{10050151}. Worked on the activation compression idea and explained quantization and sparsification. \\
\hline

Jiale Liu & Synchronization Strategy &
Reviewed papers \cite{2024ITWC...23.7362C} and \cite{2021arXiv210709786C}. Explained the synchronization interval ($\rho$) and how reducing communication rounds can improve efficiency. \\
\hline

Zhuolun Li & Scheduler and Profiler &
Reviewed papers \cite{10.3389/frwa.2024.1378598} and \cite{CHAN2023100375}. Designed and documented the Scheduler and Client Profiler modules, including the rule-based adaptation policy and per-round latency and upload-volume recording mechanism. \\
\hline

Wenjie Ding & High-level Design \& Paper Structure &
Reviewed papers \cite{2023arXiv230205599M} and \cite{2026ITWC...25.1981L}. Led the writing of the High-level Design section and was responsible for the overall logical structure and coherence of the paper. \\
\hline

\end{tabular}
\end{center}

\end{document}
